\section{Capacity is a constraint on an inference procedure}\label{sec:C11:capacity}
A context window is an admissible input length, whereas long-context intelligence is a property of a model interacting with a distribution of tasks. The difference becomes operational when an agent must combine an early instruction, a later correction, an external observation, and an unfinished commitment. A successful retrieval of an isolated sentence does not establish that the same model can reconcile those four objects. Conversely, an agent with a modest window may solve a long task by maintaining sufficient state outside that window. Input capacity, effective evidence use, and duration of operation therefore require different measurements.

Let the complete observable history be $H_t=(e_1,\ldots,e_t)$, the query or current decision request be $q_t$, and the model input be $\mathcal C_t=\phi(H_t,q_t)$. The transformation $\phi$ includes serialization, retrieval, truncation, and any instruction or tool-description overhead. A model with token budget $W$ requires $|\mathcal C_t|\leq W$, but this inequality says nothing about whether $\mathcal C_t$ contains the evidence necessary for a correct action. Define a task-dependent evidence set $D_t\subseteq H_t$ and an action loss $\ell(a,Y_t)$. The relevant risk is
\begin{equation}
 R(\phi,\pi)=\mathbb E\!\left[\ell\{\pi(\phi(H_t,q_t)),Y_t\}\right],
 \qquad |\phi(H_t,q_t)|\leq W.
\end{equation}
Here $W$ is the admitted input budget. If the serving interface constrains input and generated tokens jointly, an output reserve must first be subtracted from that joint limit. The distribution in this expectation must identify histories, decisions, environments, and evaluation times. Changing any of them can change the apparent effective context length without changing the architecture.

It is useful to distinguish a capacity failure from an accessibility failure. A capacity failure occurs when a chosen representation cannot fit a sufficient record into the available input. An accessibility failure occurs when a sufficient record is present but the model does not use it correctly. A third category, representation failure, occurs when information was transformed into an ambiguous or incorrect record before the model saw it. These categories imply different interventions: extend the budget, improve the model's use of evidence, or repair the state-construction procedure. Treating them as one phenomenon encourages expensive remedies that miss the responsible subsystem.

The transformer construction makes every permitted earlier token potentially available through attention, but potential connectivity is not reliable semantic access~\cite{vaswani2017attention}. Position encodings, training distributions, intermediate representations, decoding, and instructions all mediate that access. A systems description should therefore specify the input length, the position of decisive evidence, the number of distracting or contradictory records, and the logical operations required to derive the action. The phrase ``supports a long context'' becomes scientifically informative only after those conditions are fixed.

